% Part: turing-machines 
% Chapter: machines-computations 
% Section: introduction

\documentclass[../../../include/open-logic-section]{subfiles}

\begin{document}

\olfileid[pt]{tur}{mac}{int} 
\olsection{Introdução}

O que significa dizer que uma função, digamos, de $\Nat$ para $\Nat$ é
\emph{computável}? Entre as primeiras respostas, e a mais conhecida,
está a de que uma função é computável se puder ser computada por uma
máquina de Turing. Essa noção foi apresentada por Alan Turing em 1936.
As máquinas de Turing são um exemplo de \emph{modelo de computação}---uma
maneira matematicamente precisa de definir a ideia de um
``procedimento computacional.'' O que exatamente isso significa é debatido, mas
há amplo consenso de que as máquinas de Turing são uma forma de especificar
procedimentos computacionais. Embora o termo ``máquina de Turing''
evoque a imagem de uma máquina física com partes móveis, estritamente
falando uma máquina de Turing é um construto puramente matemático e, como
tal, idealiza a ideia de um procedimento computacional. Por
exemplo, não impomos restrição alguma aos requisitos de tempo ou memória
de uma máquina de Turing: máquinas de Turing podem computar
algo mesmo que a computação exija mais espaço de armazenamento ou
mais passos do que há átomos no universo.

\begin{explain}
Talvez seja melhor pensar em uma máquina de Turing
como um programa para um tipo especial de mecanismo imaginário. Esse
mecanismo consiste em uma \emph{fita} e uma \emph{cabeça de leitura-escrita}. Na
nossa versão das máquinas de Turing, a fita é infinita em uma direção
(à direita) e é dividida em \emph{quadrados}, cada um dos quais
pode conter um símbolo de um \emph{alfabeto} finito. Esses alfabetos podem
conter qualquer número de símbolos diferentes, mas nos contentaremos em geral
com três: $\TMendtape$, $\TMblank$ e $\TMstroke$. Quando o
mecanismo é iniciado, a fita está vazia (ou seja, cada quadrado contém
o símbolo $\TMblank$), exceto o quadrado mais à esquerda, que contém
$\TMendtape$, e um número finito de quadrados que contêm a
\emph{entrada}. A qualquer momento, o mecanismo está em um de um número finito de
\emph{estados}. No início, a cabeça examina o quadrado mais à esquerda
e em um \emph{estado inicial} especificado. A cada passo da
execução do mecanismo, o conteúdo do quadrado atualmente examinado, junto
com o estado em que o mecanismo se encontra e o programa da máquina de Turing,
determinam o que acontece em seguida. O programa da máquina de Turing é dado por uma
função parcial que recebe como entrada um estado~$q$ e um
símbolo~$\sigma$ e produz uma tripla~$\tuple{q', \sigma',
  D}$. Sempre que o mecanismo está no estado $q$ e lê o símbolo
$\sigma$, substitui o símbolo no quadrado atual por $\sigma'$, a cabeça move-se à esquerda, à direita ou permanece no lugar conforme $D$ seja
$\TMleft$, $\TMright$ ou $\TMstay$, e o mecanismo passa ao
estado~$q'$.

Por exemplo, considere a situação em \olref{fig:tm}.
\begin{figure}
  \olasset{\olpath/assets/diagrams/turing-machine.tikz}
  \caption{Uma máquina de Turing executando seu programa.}
  \ollabel{fig:tm}
\end{figure}
A parte visível da fita da máquina de Turing contém o
símbolo de fim de fita $\TMendtape$ no quadrado mais à esquerda, seguido de
três $1$'s, um $0$ e mais quatro $1$'s. A
cabeça está lendo o terceiro quadrado da esquerda, que contém
um~$1$, e está no estado~$q_1$---dizemos que ``a máquina está lendo
um $1$ no estado~$q_1$.'' Se o programa da máquina de Turing
retornar, para a entrada $\tuple{q_1, 1}$, a tripla $\tuple{q_2,
0, \TMstay}$, a máquina substituiria o~$1$ no
terceiro quadrado por um~$0$, deixaria a cabeça de leitura/escrita onde está
e passaria ao estado~$q_2$. Se então o programa retornar
$\tuple{q_3, 0, \TMright}$ para a entrada $\tuple{q_2, 0}$,
a máquina sobrescreveria o~$0$ com outro~$0$
(efetivamente, deixando inalterado o conteúdo da fita sob a cabeça de leitura/escrita), moveria um quadrado à direita e entraria no estado~$q_3$.
E assim por diante.

Dizemos que a máquina \emph{para} quando encontra algum estado,
$q_n$, e símbolo, $\sigma$, tais que não há instrução para
$\tuple{q_n, \sigma}$, ou seja, a função de transição para a entrada
$\tuple{q_n,\sigma}$ é indefinida. Em outras palavras, a máquina não tem
instrução a executar e, nesse ponto, cessa a
operação. A parada às vezes é representada por um estado de parada
específico~$h$. Isso será demonstrado com mais detalhes adiante.
\end{explain}

\begin{digress}
A beleza do artigo de Turing, ``On computable numbers,'' está em que ele
apresenta não só uma definição formal, mas também um argumento de que a
definição captura a noção intuitiva de computabilidade.
Pela definição, deve ficar claro que toda função computável
por uma máquina de Turing é computável no sentido intuitivo. Turing
oferece três tipos de argumento de que o inverso é verdadeiro, ou seja, que
toda função que naturalmente consideraríamos computável é
computável por tal máquina. São eles (nas palavras de Turing):
\begin{enumerate}
\item Um apelo direto à intuição.
\item Uma prova da equivalência de duas definições (caso a nova
  definição tenha maior apelo intuitivo).
\item Dar exemplos de grandes classes de números que são
  computáveis.
\end{enumerate}
Nosso objetivo é tentar definir a noção de computabilidade ``em
princípio,'' ou seja, sem levar em conta limitações práticas
de tempo e espaço. Claro, com a definição mais ampla de
computabilidade em vigor, pode-se então considerar computação
com recursos limitados; isso forma o cerne do assunto conhecido como
``complexidade computacional.''
\end{digress}

\begin{history}
Alan Turing inventou as máquinas de Turing em 1936. Embora seu interesse na
época fosse a decidibilidade da lógica de primeira ordem, o artigo tem sido
descrito como um trabalho definitivo sobre os fundamentos do projeto de
computadores. No artigo, Turing concentra-se em números reais computáveis, ou seja,
números reais cujas expansões decimais são computáveis; mas observa
que não é difícil adaptar suas noções a funções computáveis sobre
os números naturais, e assim por diante. Note que isso foi cinco
anos inteiros antes de o primeiro computador de propósito geral em funcionamento ser construído em
1941 (pelo alemão Konrad Zuse na sala de estar dos pais), sete
anos antes de Turing e seus colegas em Bletchley Park construírem o
Colossus que quebrava códigos (1943), nove anos antes do ENIAC americano
(1945), doze anos antes do primeiro computador britânico de propósito geral---a Manchester Small-Scale Experimental Machine---ser construída em
Manchester (1948), e treze anos antes dos americanos testarem o
BINAC (1949). A Manchester SSEM tem a distinção de
ser o primeiro computador com programa armazenado---as máquinas anteriores tinham de ser
recabeada manualmente para cada nova tarefa.
\end{history}

\end{document}
